\everymath{\displaystyle}

%%%%%%%%%%%%%%%%%%%%%%%%%%%%%%%%%%%%%%%%
%           main body
%%%%%%%%%%%%%%%%%%%%%%%%%%%%%%%%%%%%%%%%

%-----------------------------------
%	TITLE PAGE
%-----------------------------------

\title[第一章\\
集与点集]{\texorpdfstring{\LARGE \bfseries 第一章\quad 集与点集}{第一章 集与点集}} % The short title appears at the bottom of every slide, the full title is only on the title page
\author[\S 1.1\\
集及其运算] % Your institution as it will appear on the bottom of every slide, maybe shorthand to save space
{\Large \bfseries \S 1.1 集及其运算}
% \author{Singular Value Decomposition} % Your name
% \date{\today} % Date can be changed to a custom date
\date{}

%------------------------------------------------

\begin{document}

%------------------------------------------------

\begin{frame}
\frametitle{分析课程体系}

\begin{center}
\begin{tikzpicture}[
    >=Latex,
    node distance=1.1cm and 0.5cm,
    main/.style={
        draw,
        rounded corners,
        thick,
        align=center,
        fill=blue!8,
        minimum width=2.3cm,
        minimum height=1.2cm
    },
    sub/.style={
        draw,
        rounded corners,
        thick,
        align=center,
        fill=green!8,
        minimum width=2.3cm,
        minimum height=0.9cm
    },
    lower/.style={
        draw,
        rounded corners,
        thick,
        align=center,
        fill=orange!10,
        minimum width=2.3cm,
        minimum height=0.8cm
    }
]

% Top node
\node[main] (analysis) {
    {\larger[1] 数学分析}\\
    (Mathematical Analysis)\\[4pt]
    {\smaller[1]基础性、一般性的分析学知识}\\
    {\smaller[1]包括极限、微分、黎曼积分等}\\
    {\smaller[1]主要研究对象是（几乎处处）连续函数}
};

% Middle layer
\node[sub, fill=red!40, below left=of analysis] (real) {
    {\larger[1] 实分析}\\
    {\smaller[1] Function of a real variable}\\
    {\smaller[1] 一元实函数更深入的知识}
};
\node[above left=of real, xshift=1.9cm, yshift=-1cm, fill=yellow!60] (course) {本课程};

\node[sub, below=of analysis] (complex) {
    {\larger[1] 复分析}\\
    {\smaller[1] Function of a complex variable}
};

\node[sub, below right=of analysis] (functional) {
    {\larger[1] 泛函分析}\\
    {\smaller[1] Function of functions}
};

% Lower layer
\node[lower, below left=of real, xshift=2.1cm] (measure) {
    {\larger[1] 测度论}\\
    {\smaller[1] 更一般空间上的}\\
    {\smaller[1] 测度、积分理论}
};

\node[lower, below right=of real, xshift=-2.1cm] (prob) {
    概率论
};

\node[lower, below=of complex] (several) {
    多复变
};

\node[lower, below=of functional] (operator) {
    算子理论
};

% Arrows
\draw[->, thick] (analysis) -- (real);
\draw[->, thick] (analysis) -- (complex);
\draw[->, thick] (analysis) -- (functional);

\draw[->, thick] (real) -- (measure);
\draw[->, thick] (real) -- (prob);

\draw[->, thick] (complex) -- (several);

\draw[->, thick] (functional) -- (operator);
\end{tikzpicture}
\end{center}

\end{frame}

%------------------------------------------------

\begin{frame}{课程内容与主线}

\startproof[bg=yellow, fg=red](课程主要目的)
扩充、推广 Riemann 积分相关理论, 建立 Lebesgue 测度与积分相关的理论体系

\pause
\vspace{1em}

\centering
\begin{tikzpicture}[
    >=Latex,
    thick,
    every node/.style={align=center},
    main/.style={
        draw,
        rounded corners,
        fill=blue!8,
        minimum width=2.8cm,
        minimum height=1.3cm
    },
    explain/.style={
        draw,
        rounded corners,
        fill=orange!10,
        text width=2.8cm,
        minimum height=1cm,
        font=\small
    }
]

% Main horizontal line
\draw[very thick, ->] (-1.6,0) -- (13.9,0);

% Main nodes
\node[main] (set) at (0.5,0) {
    {\larger[1] 集合论}
};

\node[main] (measurable) at (4.1,0) {
    {\larger[1] 可测集与}\\
    {\larger[1] 可测函数}
};

\node[main] (integral) at (7.7,0) {
    {\larger[1] 勒贝格积分}
};

\node[main] (lp) at (11.3,0) {
    {\larger[1] $L^p$ 空间}
};

% Explanation nodes
\node[explain] (setexp) at (0.5,-2.0) {
    基础语言与工具
};

\node[explain] (meaexp) at (4.1,1.9) {
    哪些对象可以被\\
    “测量”
};

\node[explain] (intexp) at (7.7,-2.0) {
    如何对更一般的\\
    函数进行积分
};

\node[explain] (lpexp) at (11.3,1.9) {
    满足一定条件的函数\\
    所构成的空间
};

% Connecting lines
\draw[dashed] (set) -- (setexp);
\draw[dashed] (measurable) -- (meaexp);
\draw[dashed] (integral) -- (intexp);
\draw[dashed] (lp) -- (lpexp);

% Bottom philosophy arrow
\draw[ultra thick, red!60!black, ->]
    (0.0,-3.1) -- (12.5,-3.1);
\node at (0.0,-3.6) {单个函数};
\node at (12.5,-3.6) {函数空间};

% Optional subtitle
\node[font=\itshape]
    at (5.8,-4.3)
{
    {从研究“一个函数”到研究“函数组成的空间”}
};
\end{tikzpicture}

\end{frame}

%------------------------------------------------

\begin{frame}

\titlepage % Print the title page as the first slide

\end{frame}

%------------------------------------------------

% \begin{frame}
% \frametitle{内容提要} % Table of contents slide, comment this block out to remove it
% \tableofcontents % Throughout your presentation, if you choose to use \section{} and \subsection{} commands, these will automatically be printed on this slide as an overview of your presentation
% \end{frame}

%-----------------------------------
%	PRESENTATION SLIDES
%-----------------------------------

%------------------------------------------------

\section{集合论基本概念}

%------------------------------------------------

\begin{frame}{回顾: 基本的集合论知识}

我们依照朴素集合论的框架来回顾一些基本概念和记号. 关于公理化集合论, 感兴趣的同学可以参考相关书籍, 例如 Zorich 《数学分析》第 1 章第 4 节.

\vspace{0.5em}

\begin{Def}[集合]
集合是一些\alert{确定的}、\alert{互异的}对象, 或者适合\alert{某种条件}的对象的\alert{全体}.
\end{Def}

\pause

\begin{block}{基本记号}
\begin{itemize}
\item $a \in A$ ($a \notin A$): 元素的归属关系
\item $A \subset B$: $A$ 是 $B$ 的\alert{子集}，即 $\forall ~ x \in A \Rightarrow x \in B$
\begin{itemize}
\item 集合相等：$A = B \iff A \subset B$ 且 $B \subset A$
\item 真子集：$A \subset B$ 且 $A \neq B$ (或者说 $\exists ~ x \in B, x \notin A$) 记为 $A \subsetneqq B$
\end{itemize}
\item $\emptyset$：空集, 不含有任何元素的集合 ($\emptyset \subset A$ 对任意集 $A$)
\item $X$：基本集 (全集)，讨论所处的范围
\end{itemize}
\end{block}

\end{frame}

%------------------------------------------------

\begin{frame}{回顾: 基本的集合论知识}

\begin{block}{集合的表示}
\begin{itemize}
\item \alert{列举法}: $\{1, 2, 3\},$ $\mathbb{N} = \{1, 2, 3, \ldots\}$
\item \alert{描述法}: $\{x \in X : P(x)\},$ $P$ 是 $x$ 的某种性质或条件. 例如 $\{x \in \mathbb{R} : x^2 < 2\}$.
\end{itemize}
\end{block}

\pause
\vspace{0.5em}

\begin{block}{几类特殊集合}
\begin{itemize}
\item \alert{集族与指标集}: 设 $I$ 是给定集合, 对于每个 $\alpha \in I$ 赋予一个集合 $A_\alpha$, 则 $\{A_\alpha\}_{\alpha \in I}$ 是以 $I$ 为指标集 (index set) 的集族 (family of sets). 指标集 $I$ 可以是有限、可列或不可列的.
\item \alert{幂集}: 集合 $A$ 的全体子集构成的集族 $\mathscr{P}(A) = \{B : B \subset A\}$ 被称作 $A$ 的幂集 (power set). $A$ 为有限集时, $|\mathscr{P}(A)| = 2^{|A|}$.
\end{itemize}
\end{block}

\end{frame}

%------------------------------------------------

\section{集的运算}

%------------------------------------------------

\begin{frame}{集的基本运算}

设 $A, B \subset X$，基本运算包括：

\pause

\begin{itemize}
\item \alert{交} (intersection) $A \cap B$, \quad \alert{并} (union) $A \cup B$, \quad \alert{补} (complement) $\mathscr{C} A = \mathscr{C}_X A = A^c = X \setminus A$
\item \alert{差} (difference) $A \setminus B = A \cap \mathscr{C}B \;(= \{x \in A : x \notin B\})$
\item \alert{对称差} (symmetric difference) $A \triangle B = (A \setminus B) \cup (B \setminus A)$
\end{itemize}

\pause
\vspace{0.5em}

推广至任意集族 $\{A_\alpha\}_{\alpha \in I}$ ($I$ 可有限、可列或不可列):
\[
\bigcup_{\alpha \in I} A_\alpha = \{x : {\color{red}\exists} ~ \alpha \in I,\ x \in A_\alpha\},
\qquad
\bigcap_{\alpha \in I} A_\alpha = \{x : {\color{red}\forall} ~ \alpha \in I,\ x \in A_\alpha\}
\]

\pause
\vspace{0.5em}

\begin{remark}
写 $A \setminus B = A \cap \mathscr{C}B$ 可将差集化为交与补, 在推导集合等式时十分有用.
\end{remark}

\pause
\vspace{0.5em}

\begin{eg}
$A = \{ 2n - 1 : n \in \mathbb{Z} \},$ $B = \{ n : n \in \mathbb{Z}, |n| \leqslant 3 \}$, 则
\[
\begin{array}{ll}
A \cap B = \{-3, -1, 1, 3\},  & A \cup B = \{2n - 1 : n \in \mathbb{Z}\} \cup \{-2, 0, 2\}, \\
A \setminus B = \{2n - 1 : n \in \mathbb{Z}, n > 2, \text{或} ~ n \leqslant -2\}, & B \setminus A = \{-2, 0, 2\}.
\end{array}
\]
\end{eg}

\end{frame}

%------------------------------------------------

\begin{frame}{集的运算法则}

\begin{thm}[交、并的基本运算法则]
设 $A, B, C$ 是任意集合, 则有如下的 (关于集合交、并) 基本运算法则:
\begin{itemize}
\item \alert{交换律}: $A \cap B = B \cap A$, $A \cup B = B \cup A$;
\item \alert{结合律}: $(A \cap B) \cap C = A \cap (B \cap C)$, $(A \cup B) \cup C = A \cup (B \cup C)$;
\item \alert{分配律}: $A \cap (B \cup C) = (A \cap B) \cup (A \cap C)$, $A \cup (B \cap C) = (A \cup B) \cap (A \cup C)$.
\end{itemize}
\end{thm}

\pause
\vspace{0.5em}

\begin{cor}
设 $\{A_\alpha\}_{\alpha \in I}$ 是任意集族, $E$ 是任意集合, 则
\[
E \cap \!\left(\bigcup_{\alpha \in I} A_\alpha\right) = \bigcup_{\alpha \in I} (E \cap A_\alpha),
\qquad
E \cup \!\left(\bigcap_{\alpha \in I} A_\alpha\right) = \bigcap_{\alpha \in I} (E \cup A_\alpha).
\]
\end{cor}

\end{frame}

%------------------------------------------------

\begin{frame}{集的运算法则}

\begin{prop}[补、差的基本运算法则]
设 $X$ 是基本集, $A, B \subset X$ 是任意集合, 则有如下的 (关于集合补、差) 基本运算法则:
\begin{itemize}
\item $A \cup \mathscr{C}A = X,$ $A \cap \mathscr{C}A = \emptyset$;
\item $\mathscr{C}(\mathscr{C}A) = A,$ $\mathscr{C} X = \emptyset,$ $\mathscr{C} \emptyset = X$;
\item $A \setminus B = A \cap \mathscr{C}B,$ $\mathscr{C}(A \cap B) = \mathscr{C}A \cup \mathscr{C}B,$ $\mathscr{C}(A \cup B) = \mathscr{C}A \cap \mathscr{C}B$;
\item $A \subset B \iff \mathscr{C}A \supset \mathscr{C}B$;
\item $A \cap B = \emptyset \iff A \subset \mathscr{C}B \iff B \subset \mathscr{C}A$.
\end{itemize}
\end{prop}

\vspace{0.5em}

以上性质可以类似地推广至任意集族 $\{A_\alpha\}_{\alpha \in I}.$

\end{frame}

%------------------------------------------------

\begin{frame}{De Morgan 律}

\begin{thm}[De Morgan 律]
设 $\{A_\alpha\}_{\alpha \in I}$ 为集族，$I$ 为\alert{任意}指标集，则
\[
\mathscr{C}\!\left(\bigcup_{\alpha \in I} A_\alpha\right) = \bigcap_{\alpha \in I} \mathscr{C} A_\alpha,
\qquad
\mathscr{C}\!\left(\bigcap_{\alpha \in I} A_\alpha\right) = \bigcup_{\alpha \in I} \mathscr{C} A_\alpha.
\]
\end{thm}

\pause

\startproof[bg=yellow, fg=red](证明要点)

\vspace{0.3em}
以第一式为例：
\begin{align*}
x \in \mathscr{C}\!\left(\bigcup_{\alpha \in I} A_\alpha\right)
& \iff x \notin A_\alpha,\ \forall ~ \alpha \in I \\
& \iff x \in \mathscr{C} A_\alpha,\ \forall ~ \alpha \in I
\iff x \in \bigcap_{\alpha \in I} \mathscr{C} A_\alpha. \quad\square
\end{align*}

\end{frame}

%------------------------------------------------

\section{集列的上、下极限}

%------------------------------------------------

\begin{frame}{集列的上极限与下极限}

\begin{Def}[上极限集与下极限集]
设 $\{E_n\}_{n=1}^{\infty}$ 为集列，定义
\[
\varlimsup_{n\to\infty} E_n = \bigcap_{k=1}^{\infty} \bigcup_{n=k}^{\infty} E_n,
\qquad
\varliminf_{n\to\infty} E_n = \bigcup_{k=1}^{\infty} \bigcap_{n=k}^{\infty} E_n.
\]
\end{Def}

\pause

\textbf{直观解释：}
\begin{itemize}
\item $\displaystyle x \in \varlimsup_{n\to\infty} E_n \iff x$ 属于\alert{无穷多个} $E_n$
\item $\displaystyle x \in \varliminf_{n\to\infty} E_n \iff x$ 属于\alert{除有限个外所有的} $E_n$，即 $\exists ~ N,\ \forall ~ n \geqslant N,\ x \in E_n$
\end{itemize}

\pause

\begin{block}{基本性质}
\begin{itemize}
\item \alert{包含关系}:  $\displaystyle \varliminf_{n\to\infty} E_n \subset \varlimsup_{n\to\infty} E_n$
\item \alert{集列收敛}: 若两者相等，称集列收敛，极限记 $\displaystyle \lim_{n\to\infty} E_n = \varlimsup_{n\to\infty} E_n = \varliminf_{n\to\infty} E_n$
% \item \alert{单调集列}必收敛: 渐张时 $\displaystyle \lim_{n\to\infty} E_n = \bigcup_{n=1}^\infty E_n$；渐缩时 $\displaystyle \lim_{n\to\infty} E_n = \bigcap_{n=1}^\infty E_n$
\end{itemize}
\end{block}

\end{frame}

%------------------------------------------------

\begin{frame}{集列上下极限 \texorpdfstring{$|$}{|} 例}

\begin{eg}
设 $E_n = \left\{ \dfrac{m}{n} : m \in \mathbb{Z} \right\} = \dfrac{1}{n}\mathbb{Z}$，则
$\displaystyle \varliminf_{n\to\infty} E_n = \mathbb{Z}$, $\displaystyle \varlimsup_{n\to\infty} E_n = \mathbb{Q}$.
\end{eg}

\pause

\startproof[bg=yellow, fg=red](证明要点)
\vspace{0.3em}

\begin{columns}
\begin{column}[T]{0.49\textwidth}

\alert{$\displaystyle \varliminf_{n\to\infty} E_n = \mathbb{Z}$}:

\medskip
$\supset$：对任意 $p \in \mathbb{Z}$，有 $p = \frac{pn}{n} \in E_n$ 对所有 $n$ 成立，故 $\mathbb{Z} \subset \bigcap_{n=1}^\infty E_n \subset \varliminf_{n\to\infty} E_n$.

\medskip
$\subset$：若 $x \in \varliminf_{n\to\infty} E_n$，则从某个 $N$ 起 $x \in E_n$. 设 $x = p/q$（既约, $q > 0$），可证必有 $q = 1$，即 $x \in \mathbb{Z}$.

\end{column}
\begin{column}[T]{0.49\textwidth}

\alert{$\displaystyle \varlimsup_{n\to\infty} E_n = \mathbb{Q}$}:

\medskip
$\subset$：由 $E_n \subset \mathbb{Q}$ 对所有 $n$ 成立，得 $\displaystyle \varlimsup_{n\to\infty} E_n \subset \mathbb{Q}$.

\medskip
$\supset$：对任意 $\frac{p}{q} \in \mathbb{Q}$ ($q > 0$, 既约) 和任意 $k \in \mathbb{N}$，取 $n = kq$，则 $\frac{p}{q} = \frac{kp}{kq} \in E_n$，故 $\frac{p}{q} \in \bigcup_{n \geqslant k} E_n$ 对所有 $k$ 成立.

\end{column}
\end{columns}

\end{frame}

%------------------------------------------------

\end{document}
